Par déplacement d'air, à $\qty{20}{\degreeCelsius}$ et sous la pression de
$\qty{101300}{\Pa}$, on introduit du chlorure d'hydrogène dans un flacon de
$\qty{500}{\mL}$ . On réalise ensuite l'expérience du jet d'eau et on constate
qu'en fin d'expérience, on recueille un volume de solution de $\qty{200}{\mL}$.
\begin{enumerate}
    \item Comment recueille-t-on un gaz par déplacement d'air~?
    \item Calculer la quantité de matière de chlorure d'hydrogène présente, au
          départ, dans le flacon.
    \item Calculer la concentration molaire de la solution d'acide chlorhydrique
          obtenue.
\end{enumerate}
\begin{donnees}
    $R=\qty{8.31}{SI}$.
\end{donnees}

